\documentclass[journal]{vgtc}                % final (journal style)
%\documentclass[review,journal]{vgtc}         % review (journal style)
%\documentclass[widereview]{vgtc}             % wide-spaced review
%\documentclass[preprint,journal]{vgtc}       % preprint (journal style)

%% Uncomment one of the lines above depending on where your paper is
%% in the conference process. ``review'' and ``widereview'' are for review
%% submission, ``preprint'' is for pre-publication, and the final version
%% doesn't use a specific qualifier.

%% Please use one of the ``review'' options in combination with the
%% assigned online id (see below) ONLY if your paper uses a double blind
%% review process. Some conferences, like IEEE Vis and InfoVis, have NOT
%% in the past.

%% Please note that the use of figures other than the optional teaser is not permitted on the first page
%% of the journal version.  Figures should begin on the second page and be
%% in CMYK or Grey scale format, otherwise, colour shifting may occur
%% during the printing process.  Papers submitted with figures other than the optional teaser on the
%% first page will be refused. Also, the teaser figure should only have the
%% width of the abstract as the template enforces it.

%% These few lines make a distinction between latex and pdflatex calls and they
%% bring in essential packages for graphics and font handling.
%% Note that due to the \DeclareGraphicsExtensions{} call it is no longer necessary
%% to provide the the path and extension of a graphics file:
%% \includegraphics{diamondrule} is completely sufficient.
%%
\ifpdf%                                % if we use pdflatex
  \pdfoutput=1\relax                   % create PDFs from pdfLaTeX
  \pdfcompresslevel=9                  % PDF Compression
  \pdfoptionpdfminorversion=7          % create PDF 1.7
  \ExecuteOptions{pdftex}
  \usepackage{graphicx}                % allow us to embed graphics files
  \DeclareGraphicsExtensions{.pdf,.png,.jpg,.jpeg} % for pdflatex we expect .pdf, .png, or .jpg files
\else%                                 % else we use pure latex
  \ExecuteOptions{dvips}
  \usepackage{graphicx}                % allow us to embed graphics files
  \DeclareGraphicsExtensions{.eps}     % for pure latex we expect eps files
\fi%

%% it is recommended to use ``\autoref{sec:bla}'' instead of ``Fig.~\ref{sec:bla}''
\graphicspath{{figures/}{pictures/}{images/}{./}} % where to search for the images

\usepackage{microtype}                 % use micro-typography (slightly more compact, better to read)
\PassOptionsToPackage{warn}{textcomp}  % to address font issues with \textrightarrow
\usepackage{textcomp}                  % use better special symbols
\usepackage{mathptmx}                  % use matching math font
\usepackage{times}                     % we use Times as the main font
\renewcommand*\ttdefault{txtt}         % a nicer typewriter font
\usepackage{cite}                      % needed to automatically sort the references
\usepackage{tabu}                      % only used for the table example
\usepackage{booktabs}                  % only used for the table example
%% We encourage the use of mathptmx for consistent usage of times font
%% throughout the proceedings. However, if you encounter conflicts
%% with other math-related packages, you may want to disable it.

%% In preprint mode you may define your own headline.
%\preprinttext{To appear in IEEE Transactions on Visualization and Computer Graphics.}

%% If you are submitting a paper to a conference for review with a double
%% blind reviewing process, please replace the value ``0'' below with your
%% OnlineID. Otherwise, you may safely leave it at ``0''.
\onlineid{0}

%% declare the category of your paper, only shown in review mode
\vgtccategory{Research}
%% please declare the paper type of your paper to help reviewers, only shown in review mode
%% choices:
%% * algorithm/technique
%% * application/design study
%% * evaluation
%% * system
%% * theory/model
\vgtcpapertype{application/design study}

%% Paper title.
\title{How climate change affects bird migration}

%% This is how authors are specified in the journal style

%% indicate IEEE Member or Student Member in form indicated below
\author{Gabriela Slavova, Henrique Dias, Nimo Beeren}
\authorfooter{
%% insert punctuation at end of each item
\item
Gabriela Slavova, Henrique Dias and Nimo Beeren are with Eindhoven University of
Technology. E-mail: \{g.slavova, h.a.coelho.dias, n.beeren\}@student.tue.nl.
}

%other entries to be set up for journal
%\shortauthortitle{Biv \MakeLowercase{\textit{et al.}}: Global Illumination for Fun and Profit}
%\shortauthortitle{Firstauthor \MakeLowercase{\textit{et al.}}: Paper Title}

%% Abstract section.
\abstract{Abstract text%
} % end of abstract

%% Keywords that describe your work. Will show as 'Index Terms' in journal
%% please capitalize first letter and insert punctuation after last keyword
\keywords{Climate change, bird migration.}

%% ACM Computing Classification System (CCS). 
%% See <http://www.acm.org/class/1998/> for details.
%% The ``\CCScat'' command takes four arguments.

%\CCScatlist{ % not used in journal version
% \CCScat{K.6.1}{Management of Computing and Information Systems}%
%{Project and People Management}{Life Cycle};
% \CCScat{K.7.m}{The Computing Profession}{Miscellaneous}{Ethics}
%}

%% Uncomment below to include a teaser figure.
%\teaser{
%  \centering
%  \includegraphics[width=\linewidth]{CypressView}
%  \caption{In the Clouds: Vancouver from Cypress Mountain. Note that the teaser may not be wider than the abstract block.}
%	\label{fig:teaser}
%}

%% Uncomment below to disable the manuscript note
%\renewcommand{\manuscriptnotetxt}{}

%% Copyright space is enabled by default as required by guidelines.
%% It is disabled by the 'review' option or via the following command:
% \nocopyrightspace

\vgtcinsertpkg

%%%%%%%%%%%%%%%%%%%%%%%%%%%%%%%%%%%%%%%%%%%%%%%%%%%%%%%%%%%%%%%%
%%%%%%%%%%%%%%%%%%%%%% START OF THE PAPER %%%%%%%%%%%%%%%%%%%%%%
%%%%%%%%%%%%%%%%%%%%%%%%%%%%%%%%%%%%%%%%%%%%%%%%%%%%%%%%%%%%%%%%

\begin{document}

%% The ``\maketitle'' command must be the first command after the
%% ``\begin{document}'' command. It prepares and prints the title block.

%% the only exception to this rule is the \firstsection command
\firstsection{Introduction}

\maketitle

In recent times, it has become increasingly clear to most people that climate change is affecting life on Earth. As temperatures rise globally, many animals are forced to adapt to their environment \cite{parmesan2007pheno,gaughan2009domestic,root2003fingerprints}. Among the best-documented instances of such adaptation are changes in bird migration times \cite{miller2008bird,visser2008climate,jenni2003timing}. This paper presents an interactive visualization with the goal to communicate these findings.

\subsection{Problem Description}

Despite evidence presented in numerous studies \cite{solomon2007climate,parry2007climate}, it has proved to be difficult to communicate the effects of climate change to people \cite{lee2015predictors,brulle2012shifting,moser2011communicating}. We believe that visualizing effects that can be observed in everyday life is an important way of fostering public engagement. One such effect is the change in the time of the year during which migratory birds can be observed in their breeding habitat. As global temperatures increased, birds have been migrating to their breeding habitats earlier in the year \cite{cotton2003avian,huppop2003north,marra2005influence}. Using a visual approach, we intend to show this correlation in a way that is interpretable to experts and non-experts alike.

Both bird migration and surface temperature changes are cyclic processes with an annual period. In order to infer trends in these processes, we should consider data from multiple years, or even decades. However, because of the cyclic nature, it is difficult to directly compare one year to another. One of the challenges of this problem is to visualize trends over the long term, without obscuring the patterns that occur every year.

Another attribute that these processes have in common is that they are both geographic. That is, both produce data that is associated with a particular place on Earth \cite{iso2014geo}. It is essential to visualize both in such a way that one does not obscure the other, while their correlation is still visible.

\section{Data Analysis}

\subsection{Domain Data Specification}

In order to visualize the problem, we need two different datasets: one that includes information about birds migration and one including temperature changes over the years.

For the first needed dataset we have chosen the eBird Basic Dataset \cite{ebird2020data}. The eBird citizen science project is unmatched in scale, containing over 600 million observations, including nearly every bird species on Earth \cite{strimas2020ebird}. Submissions are individually reviewed to ensure accuracy of information. As people are contributing to this dataset continuously, the information inside is up-to-date, but the record also goes many years back.

For the second, we have chosen the Terraclimate dataset \cite{terraclimatedata} by Climatology Lab, which provides diverse information about climate - precipitation, minimum and maximum temperatures, vapor pressure, etc - over the years, since 1958. For the use case in study, we only use the maximum temperature values. One of the main advantages of this dataset is the use of climatically aided interpolation, combining values from the WorldClim \cite{worldclim} dataset with other datasets such as CRU Ts4.0 \cite{cru} and JRA55 \cite{j55r}.

\subsection{Data Abstraction: What}

The eBird Basic Dataset \cite{ebird2020data} consists of a large static table containing items that each correspond to a particular observation of a bird or a group of birds. For each observation, a number of attributes are provided, though we are only interested in the following subset:

\paragraph{Species} A categorical attribute represented by the scientific name, e.g. \textit{Anas platyrhynchos}. The eBird taxonomy\cite{ebird2019taxonomy} specifies exactly what is meant with these names.

\paragraph{Date} An ordered, quantitative and diverging attribute in year-month-day format (yyyy-mm-dd).

\paragraph{Location} A spatial attribute consisting of a vector with components latitude and longitude. Each component is itself an ordered, quantitive and cyclic attribute represented as decimal degrees, where positive values denote N/W and negative values denote S/E for latitude/longitude respectively. The precision varies between observations, but is typically 0.001 degrees or more precise (which is less than 1 km on Earth's surface).

\paragraph{Count} An ordered, quantitative and diverging attribute represented as an integer. Since participants are allowed to group multiple birds in a single observation, the value can be 1 or greater.

\vspace{2mm}

We perform a number of transformations on the original dataset for the purpose of our visualization. First, we limit ourselves to a single species, since different species exhibit significantly different migratory behavior. We chose the Barn Swallow (\textit{Hirundo rustica}) because it is a relatively common migratory bird, and it typically migrates over a large distance \cite{turner1989swallow}. Next, we restrict the dataset geographically to Europe and Africa. These transformations allow us to focus on a general population of birds. Finally, we exclude observations outside the years 1950-2020, since that matches the domain of the temperature dataset we use.

\vspace{2mm}

The Terraclimate dataset \cite{terraclimatedata} consists of a large static muldimensional field containing items where each corresponds to a certain observation in a certain location and point in time. For each observation, a few attributes are provided, though we are only interested in the following subset:

\paragraph{Latitude} An ordered quantitative cyclic attribute represented as decimal degrees, varying from -90.0 to 90.0.

\paragraph{Longitude} An ordered quantitative cyclic attribute represented as decimal degrees, varying from -180.0 to 180.0.

\paragraph{Maximum temperature} A quantitive attribute represented as decimal degrees, in Celcius degrees.

\vspace{2mm}

Although other variables are available, such as the minimum temperature and precipitation, we chose to use the maximum temperature. By doing so, we can better visualize how the increasing temperatures have been influencing the dates of bird migrations.

% TODO: mention transformations
% note that the 2020 file might be released quite soon!

\section{Task Analysis}

\subsection{Domain Specific Tasks}
% In this section, you must analyze the possible visualization tasks that an expert in the chosen domain would be interested in. Ideally, you would present a set of tasks and questions that are intimately related to the chosen domain.
This visualization is designed to be used by a large variety of specialists and non-specialists alike. An ornithologist could study the migraiton of the Barn Swallow using the provided visualization. It is also suitabale for the needs of climate  change activists who wants prove a point about a real prolem or a curious person without specific domain knowledge who wants to see the effects of climate change on the nature. The tasks that  these people could undertake a described below.

\paragraph{Task 1: Discover correlation between temperature change and the timing of bird migration} One of the thing that can be achieve with this visualizaiton is showing the migraiton of birds with relation to the change in temperature in the past few years. This can be done by choosing a bird species and comparing timing of their migraiton with respect to the temperatures for the specific time and the same migration in a different time span before or after the chosen one.  Question that would be answered by this task is: What is the correlation between temparature change and the timing of bird migration?

\paragraph{Task 2: Characterize the migration behavior of a species} Another thing that is possible to be visualized is the migration pattern of a specific type of bird over the course of a year. This can be done be choosing a bird species and following its migration behavior through a chosen year. Question that would be answered by this task is: What is the general migration behavior pattern for bird species X?

 \paragraph{Task 3: Identify species that respond most to temperature change} Diffrent types of birds respond defferently to a change in temperatures. The tool can help visualize which types of birds are affected by the rising temperatures more than others  This can be done by choosing all available types of birds or some specificaly chosen ones and comparing their migraiton behavior over the span of couple of years with relation to the temperatures during that time. Question that would be answered by this task is: Which types of birds respond to temperature change the most?


\subsection{Task Abstraction: Why}

% Once you have defined a set of domain specific tasks, you must generalize them to be domain independent to be able to use the principles from visualization. In this sense, we propose you follow the principles described in Munzner's book \cite{munzner2014visualization} and described in the lectures.

\section{Visualization and Interaction Design: How}

% Describe and motivate your choices for the visual encodings and interaction design (\emph{How}). You should not think about implementation yet, but rather what visual encodings fit best to address the task/data results presented earlier. You must justify your choices based on the design principles described in the course. Here is where you explicitly link your tasks, questions and data to the design, and argue why the chosen design will enable users to perform the tasks and answer their questions. Explain the main interactions provided in your solution.

\section{Realization}

% After you have a visualization design, you should elaborate on what you have implemented and what compromises were made. For instance, your ideal solution may consider edge bundling \cite{holten2006hierarchical} but you could not implement it in your final solution. In this section, you should also mention what framework/technology you used for your implementation and justify why you chose it.

\section{Use Cases}

% Describe the analysis of the data and how your application was used to perform the analysis. Here you demonstrate how to use your application to perform the tasks and answer the questions as an evaluation of your design and implementation.

\section{Discussion and Conclusion}

% Finally, you will reflect on what you have achieved. For instance, this section should discuss whether you managed to do what you initially planned, whether your initial choices worked well or not, things that you discovered that were not correct, etc. Conclusions about your whole work would be provided.

\bibliographystyle{abbrv}
\bibliography{main}
\end{document}

